% TOP_PROBLEM: {"id": 3316, "title": "Degenerate colorings of planar graphs", "category_id": 3, "category": {"id": 3, "name": "graph_theory", "display_name": "Graph Theory", "description": "Problems involving graphs, networks, and their properties.", "slug": "graph-theory", "order_index": 3, "created_at": "2026-07-31T15:26:25.671Z"}, "status": "open", "problem_number": "OPG-798", "set_id": 12}
% FIRST_POSTED: 2026-10-01
% ENTRY_KIND: statement_only
% UnsolvedMath ID 3316; source catalogue code OPG-798
% !TeX program = lualatex
% Definition and sources only; this file is not an LLM solution attempt.
\documentclass[11pt,a4paper]{article}
\usepackage[margin=25mm]{geometry}
\usepackage{fontspec}
\setmainfont{lmroman10-regular.otf}[BoldFont=lmroman10-bold.otf,
 ItalicFont=lmroman10-italic.otf,BoldItalicFont=lmroman10-bolditalic.otf]
\usepackage{amsmath,amssymb,mathtools}
\usepackage{unicode-math}
\setmathfont{latinmodern-math.otf}
\AtBeginDocument{\let\setminus\smallsetminus}
\usepackage{xurl,hyperref}
\hypersetup{colorlinks=true,urlcolor=blue}
\setcounter{secnumdepth}{0}
\setlength{\parindent}{0pt}
\setlength{\parskip}{5pt}
\setlength{\emergencystretch}{3em}
\begin{document}
\section*{Degenerate colorings of planar graphs}\label{3316}
{\small UnsolvedMath ID 3316 · OPG-798 · Graph Theory · OpenGarden}
\subsection{Definitions and mathematical statement}
A finite graph $H$ is $d$-degenerate if every nonempty subgraph of $H$
has a vertex of degree at most $d$, with degree measured inside that
subgraph. Equivalently, its vertices admit an ordering in which each
vertex has at most $d$ neighbours later in the ordering. In particular,
0-degenerate graphs are edgeless and 1-degenerate graphs are forests.

Let $G$ be a finite simple planar graph. The conjecture asks for a
colouring $c:V(G)\to\{1,2,3,4,5\}$ such that, for every
$k\in\{1,2,3,4\}$ and every set $I\subseteq\{1,2,3,4,5\}$ of
$k$ colour names,
\[
                         G[c^{-1}(I)]
                       \text{ is }(k-1)\text{-degenerate}.
\]
Unused colours and empty induced subgraphs are allowed. The $k=1$
condition ensures the colouring is proper. The $k=2$ condition then
forbids all two-coloured cycles; the $k=3$ and $k=4$ conditions demand
additional deletion orderings for the corresponding induced graphs.
The ordering may depend on the chosen colour subset $I$.

The range stops at four. Requiring the same property for all five
colours would force the whole graph to be 4-degenerate, which is
false for planar graphs of minimum degree five. Thus the record's
particular five-colour statement should not be strengthened by
extending the displayed quantifier to $k=5$. The problem asks that
all specified colour-subset conditions hold for one common colouring,
not that each condition can be met by a different colouring.
\subsection{Short English statement}
Can every planar graph receive five colours so that any one, two, three, or four colour classes induce graphs of degeneracy at most zero, one, two, or three respectively?
\subsection{Sources}
\begin{itemize}
\item[S1] ULAM.AI and the UnsolvedMath Contributors, supplied problems.json, record 3316 (OPG-798), OpenGarden collection. Catalogue identity and source transcription; CC BY 4.0.
\url{https://huggingface.co/datasets/ulamai/UnsolvedMath}.
\item[S2] Open Problem Garden, Degenerate colorings of planar graphs. Original problem and discussion; the mathematical formulation below is expanded from the supplied source record.
\url{http://www.openproblemgarden.org/op/degenerate_colorings_of_planar_graphs}.
\end{itemize}
\end{document}
